\section{Solutions}\label{sec:solutions_fields}

% Solution by Gemini 3.7 Flash (High)
\begin{exercisesolution}[Solution to \cref{exc:finite_field_f5_arithmetic}]
\label{sol:finite_field_f5_arithmetic}
\exinfo{Official Solution (Problem Sheet 3, Exercise 2). Performs arithmetic calculations in the finite field $\mathbb{F}_5$.}
In $\mathbb{F}_5 = \{0, 1, 2, 3, 4\}$, arithmetic is performed modulo $5$.
\begin{enumerate}[label=\textbf{(\alph*)}]
    \item We seek all pairs $(x, y) \in \mathbb{F}_5 \times \mathbb{F}_5$ such that $x + y = 0 \pmod 5$, which means $y = -x$:
    \[
    (0, 0), \quad (1, 4), \quad (2, 3), \quad (3, 2), \quad (4, 1).
    \]

    \item We compute $3^4$ and $3^{-1} = \frac{1}{3}$ in $\mathbb{F}_5$:
    \[
    3^2 = 9 \equiv 4 \pmod 5 \implies 3^4 = (3^2)^2 \equiv 4^2 = 16 \equiv 1 \pmod 5.
    \]
    For the multiplicative inverse, $3 \cdot 2 = 6 \equiv 1 \pmod 5$, so $\frac{1}{3} = 3^{-1} = 2$. Therefore:
    \[
    3^4 + \frac{1}{3} = 1 + 2 = 3 \in \mathbb{F}_5.
    \]

    \item Notice that $4 \equiv -1 \pmod 5$. Therefore:
    \[
    4^{2025} \equiv (-1)^{2025} = -1 \equiv 4 \pmod 5.
    \]
    Thus, $4^{2025} = 4 \in \mathbb{F}_5$.
\end{enumerate}
\end{exercisesolution}

% Solution by Gemini 3.7 Flash (High)
\begin{exercisesolution}[Solution to \cref{exc:frobenius_f3}]
\label{sol:frobenius_f3}
\exinfo{Official Solution (Problem Sheet 3, Exercise 3). Expands $(a+b)^3 = a^3 + b^3$ in characteristic 3 via the binomial formula.}
Let $a, b \in \mathbb{F}_3$. Using the binomial expansion formula (which holds in any commutative ring):
\[
(a + b)^3 = a^3 + 3a^2b + 3ab^2 + b^3.
\]
In the field $\mathbb{F}_3$, the integer $3$ is equivalent to $0$ ($3 \cdot 1 = 1 + 1 + 1 = 0$). Therefore:
\[
3a^2b = 0 \cdot (a^2b) = 0 \quad \text{and} \quad 3ab^2 = 0 \cdot (ab^2) = 0.
\]
Substituting these values back into the expansion gives:
\[
(a + b)^3 = a^3 + 0 + 0 + b^3 = a^3 + b^3.
\]
This proves that the Frobenius map $x \mapsto x^3$ is an additive endomorphism of $\mathbb{F}_3$.
\end{exercisesolution}

% Solution by Gemini 3.7 Flash (High)
\begin{exercisesolution}[Solution to \cref{exc:quadratic_field_extensions}]
\label{sol:quadratic_field_extensions}
\exinfo{Official Solution (Problem Sheet 3, Exercise 5). Verifies the field axioms for the quadratic field extension $k[\tau]$ using the norm function.}
\begin{enumerate}[label=\textbf{(\alph*)}]
    \item We verify the field axioms for $(k[\tau], +, \cdot, 0, 1)$:
    \begin{itemize}
        \item \textbf{Axioms (K1)--(K4) [Addition]:} Addition is defined componentwise as $(a+b\tau) + (c+d\tau) = (a+c) + (b+d)\tau$. Since $(k, +)$ is an abelian group, addition on $k[\tau]$ is associative and commutative, $0 = 0 + 0\tau$ is the additive identity, and $-(a+b\tau) = (-a) + (-b)\tau$ is the additive inverse.
        
        \item \textbf{Axioms (K5)--(K8) [Multiplication and Distributivity]:}
        Multiplication is commutative:
        \[
        (a+b\tau)(c+d\tau) = (ac + \alpha bd) + (bc + ad)\tau = (ca + \alpha db) + (da + cb)\tau = (c+d\tau)(a+b\tau).
        \]
        The element $1 = 1 + 0\tau$ satisfies $(a+b\tau)(1+0\tau) = (a\cdot 1 + \alpha b\cdot 0) + (b\cdot 1 + a\cdot 0)\tau = a+b\tau$, so $1$ is the multiplicative identity. Associativity and distributivity follow by direct algebraic expansion from the corresponding properties in $k$ together with the relation $\tau^2 = \alpha$.
        
        \item \textbf{Axiom (K7) [Multiplicative Inverse]:} Let $z = a + b\tau \in k[\tau]$ with $z \ne 0$, meaning $(a, b) \ne (0, 0)$. Define the \newterm{norm} of $z$ by:
        \[
        N(z) := (a + b\tau)(a - b\tau) = a^2 - \alpha b^2 \in k.
        \]
        We claim that $N(z) \ne 0$ in $k$. Indeed, if $N(z) = 0$, then $a^2 = \alpha b^2$.
        \begin{itemize}
            \item If $b \ne 0$, we can multiply by $(b^{-1})^2$ to obtain $(ab^{-1})^2 = \alpha$, which means $x = ab^{-1} \in k$ is a solution to $x^2 = \alpha$, contradicting the hypothesis that $x^2 = \alpha$ has no solution in $k$.
            \item If $b = 0$, then $a^2 = 0 \implies a = 0$, which contradicts $(a, b) \ne (0, 0)$.
        \end{itemize}
        Thus $N(z) = a^2 - \alpha b^2 \ne 0$ in $k$. Since $k$ is a field, $N(z)$ has a multiplicative inverse $N(z)^{-1} \in k$. Setting
        \[
        z^{-1} := \frac{a}{a^2 - \alpha b^2} - \frac{b}{a^2 - \alpha b^2}\tau \in k[\tau],
        \]
        we compute $z \cdot z^{-1} = \frac{(a+b\tau)(a-b\tau)}{a^2 - \alpha b^2} = \frac{a^2 - \alpha b^2}{a^2 - \alpha b^2} = 1$. Thus, every non-zero element has a multiplicative inverse.
    \end{itemize}
    This completes the proof that $k[\tau]$ is a field.

    \item \textbf{Examples:}
    \begin{itemize}
        \item \textbf{The Complex Numbers:} Take $k = \mathbb{R}$ and $\alpha = -1$. The equation $x^2 = -1$ has no solution in $\mathbb{R}$. Adjoining $\tau = i$ gives $\mathbb{R}[i] \cong \mathbb{C}$.
        \item \textbf{Quadratic Number Fields:} Take $k = \mathbb{Q}$ and $\alpha = 2$. Since $\sqrt{2} \notin \mathbb{Q}$, $\mathbb{Q}[\sqrt{2}] = \{a + b\sqrt{2} \mid a, b \in \mathbb{Q}\}$ is a field extension of $\mathbb{Q}$.
        \item \textbf{Finite Fields:} Take $k = \mathbb{F}_3 = \{0, 1, 2\}$ and $\alpha = 2 \equiv -1$. The squares in $\mathbb{F}_3$ are $0^2 = 0, 1^2 = 1, 2^2 = 1$, so $x^2 = 2$ has no solution in $\mathbb{F}_3$. Then $\mathbb{F}_3[\tau]$ is a finite field with $3^2 = 9$ elements, isomorphic to $\mathbb{F}_9$.
    \end{itemize}
\end{enumerate}
\end{exercisesolution}

% Solution by Gemini 3.7 Flash (High)
\begin{exercisesolution}[Solution to \cref{exc:finite_field_sum_product}]
\label{sol:finite_field_sum_product}
\exinfo{Official Solution (Problem Sheet 3, Exercise 6). Proves the sum $\sum_{x\in k} x = 0$ (for $|k|>2$) via scaling permutation and Wilson's theorem for finite fields $\prod_{x\in k^\times} x = -1$.}
\begin{enumerate}[label=\textbf{(\alph*)}]
    \item \textbf{Sum of all elements:}
    \begin{itemize}
        \item Suppose $|k| > 2$. Then the multiplicative group $k^\times = k \setminus \{0\}$ has cardinality $|k^\times| = |k| - 1 \ge 2$, so there exists an element $b \in k^\times$ with $b \ne 1$.
        The multiplication map $m_b \colon k \to k, x \mapsto bx$ is a bijection (with inverse $m_{b^{-1}}$). Therefore, multiplying every term in the sum $S := \sum_{x \in k} x$ by $b$ simply permutes the elements of $k$:
        \[
        bS = b \sum_{x \in k} x = \sum_{x \in k} bx = \sum_{y \in k} y = S.
        \]
        This implies $(1 - b)S = 0$. Since $b \ne 1$, $1 - b \ne 0$, and by \cref{lem:no_zero_divisors} (absence of zero divisors in a field), we conclude that $S = 0$.
        \item Conversely, if $|k| \le 2$, since any field must have distinct elements $0 \ne 1$ (by axiom \textbf{K10}), the only possibility is $k = \mathbb{F}_2 = \{0, 1\}$. In this case, $S = 0 + 1 = 1 \ne 0$.
    \end{itemize}
    Hence, $S = 0$ if and only if $|k| > 2$.

    \item \textbf{Product of all non-zero elements:} Consider the inversion map $\iota \colon k^\times \to k^\times, x \mapsto x^{-1}$. Since inversion is a bijection from $k^\times$ to itself, every element $x$ can be paired with its unique inverse $x^{-1}$, and their product is $x \cdot x^{-1} = 1$.
    
    An element is its own inverse if and only if
    \[
    x = x^{-1} \iff x^2 = 1 \iff x^2 - 1 = 0 \iff (x - 1)(x + 1) = 0.
    \]
    By \cref{lem:no_zero_divisors}, the only solutions in $k$ are $x = 1$ and $x = -1$.
    \begin{itemize}
        \item If $\operatorname{char}(k) \ne 2$, then $1 \ne -1$, so exactly two elements ($1$ and $-1$) are self-inverse. All remaining elements of $k^\times \setminus \{1, -1\}$ pair up into disjoint two-element sets $\{y, y^{-1}\}$ with $y \cdot y^{-1} = 1$. Therefore:
        \[
        M = \prod_{x \in k^\times} x = 1 \cdot (-1) \cdot \prod_{\{y, y^{-1}\}} (y \cdot y^{-1}) = -1 \cdot 1 = -1.
        \]
        \item If $\operatorname{char}(k) = 2$, then $1 = -1$, so $x = 1$ is the unique self-inverse element. All other elements pair up as $\{y, y^{-1}\}$ with product $1$. Therefore:
        \[
        M = 1 \cdot \prod_{\{y, y^{-1}\}} 1 = 1 = -1 \quad (\text{since } 1 = -1 \text{ in characteristic } 2).
        \]
    \end{itemize}
    In both cases, we have $M = -1$.
\end{enumerate}
\end{exercisesolution}
